% !TEX root = ../main.tex

\chapter{绪论}

\section{研究背景与意义}

近年来，大语言模型在自然语言理解、文本生成、逻辑推理和代码生成等任务中取得了显著进展。以 Transformer 为代表的神经网络结构通过自注意力机制增强了模型对长距离依赖关系的建模能力，为后续大规模预训练语言模型的发展奠定了基础\cite{Vaswani2017Attention}。在此基础上，研究者开始将大规模预训练方法应用于程序语言建模，使模型能够从海量代码语料中学习编程语言的语法结构、常用 API 调用模式和问题求解范式。与普通自然语言相比，程序语言具有更强的形式化结构和更明确的可执行语义，因此代码生成任务既受到自然语言生成技术发展的推动，也具有区别于普通文本生成的特殊研究价值。

代码生成任务通常要求模型根据自然语言描述、函数签名或上下文代码自动生成满足需求的程序片段。该任务能够应用于自动化编程、智能代码补全、程序修复、测试用例生成和软件工程教育等场景。随着 Codex、Code Llama、DeepSeek-Coder 等代码大模型的发展，模型在 HumanEval、MBPP 等基准测试上的表现不断提升，表明大规模预训练模型已经具备一定的程序合成能力\cite{Chen2021Codex,Roziere2023CodeLlama,Guo2024DeepSeekCoder}。然而，代码生成模型的输出并不总是可靠。自然语言生成任务中可接受的表达差异，在代码任务中可能直接表现为语法错误、运行时异常、入口函数不匹配或边界条件处理错误。即使生成代码在表面上符合题意，也可能无法通过单元测试，因而代码生成任务更强调功能正确性和可执行性。

目前，代码大模型常通过监督微调进一步适配特定任务格式。监督微调能够使模型学习给定数据集中的 prompt-response 格式，并提升模型对任务指令的响应一致性。指令微调和人类反馈训练的成功表明，预训练模型可以通过后训练阶段更好地对齐用户意图\cite{Ouyang2022InstructGPT}。但对于代码生成任务而言，仅依靠标准答案进行监督学习存在明显局限。一方面，代码问题往往存在多种正确实现，模型模仿某一个参考答案并不等价于学会满足测试语义；另一方面，监督微调目标通常基于 token 级别似然优化，而代码质量的关键评价标准是整个程序是否可执行、是否符合函数接口、是否通过测试。也就是说，训练目标与最终评价目标之间存在不一致。

基于执行反馈的强化学习为缓解上述问题提供了可行思路。代码具有天然的可执行性，模型生成的程序可以被解释器、编译器或测试框架直接验证。与人工偏好反馈相比，执行结果能够提供更客观的任务反馈。例如，代码是否可以运行、是否出现异常、通过了多少测试用例、是否全部通过测试，都可以被转化为奖励信号。CodeRL、PPOCoder 和 COMPCODER 等研究已经表明，将编译反馈、单元测试反馈或结构对齐信息引入训练过程，有助于提升模型生成代码的可执行性和功能正确性\cite{Le2022CodeRL,Shojaee2023PPOCoder,Wang2022CompCoder}。因此，围绕执行反馈设计强化学习训练流程，是提升代码生成模型可靠性的一个重要方向。

本课题以“基于执行反馈的强化学习代码生成研究”为题，关注在小规模代码数据条件下，如何构建从监督微调到强化学习优化的完整实验流程。具体而言，本文以 Qwen2.5-Coder-1.5B 作为基座模型，基于 MBPP 和 HumanEval 数据集构建统一格式的数据文件，先通过监督微调获得任务适配模型，再利用模型生成代码后的执行结果构造奖励，并采用 PPO 算法对策略进行进一步优化。本文的研究意义主要体现在以下几个方面。

第一，从训练目标角度看，本文尝试将代码任务的功能正确性纳入模型优化过程。相比仅优化参考答案似然，执行反馈能够直接反映生成程序是否满足测试约束，更接近代码生成任务的最终目标。第二，从工程实现角度看，本文构建了包括数据处理、SFT 训练、代码生成、执行打分、奖励计算、优势估计和 PPO 更新在内的完整闭环，为后续扩展更复杂的代码强化学习方法提供了实验基础。第三，从实验分析角度看，本文不仅比较 Base、SFT 和 PPO 的最终结果，还进一步分析不同 reward 设计、学习率、KL 系数和 rollout 数量对结果的影响，从而更清楚地揭示执行反馈强化学习在小规模代码数据条件下的有效性与局限性。

\section{国内外研究现状}

代码生成研究的发展大致经历了从规则驱动、统计学习、神经网络生成到大语言模型生成的过程。早期程序合成方法更多依赖形式化约束、搜索算法和手工设计的程序空间，而近年来的研究逐渐转向利用大规模神经网络从自然语言和代码语料中学习生成模式。随着大语言模型规模和数据规模不断扩大，代码生成模型已经能够在大量常见编程任务上生成可运行代码，但其功能正确性、鲁棒性和对复杂需求的理解能力仍然存在不足。

\subsection{代码生成模型研究现状}

预训练代码大模型是当前代码生成研究的核心方向之一。Codex 系列工作系统评估了大语言模型在代码生成任务中的能力，并提出 HumanEval 作为衡量函数级程序合成能力的重要基准\cite{Chen2021Codex}。Austin 等人进一步探讨了大语言模型在程序合成任务中的表现，说明大模型可以通过从自然语言描述中学习程序模式完成一定复杂度的代码生成任务\cite{Austin2021ProgramSynthesis}。这些研究推动了代码生成从传统程序合成方法向大规模预训练模型方法转变。

随着开源代码大模型的发展，Code Llama 和 DeepSeek-Coder 等模型进一步提升了开放代码模型的能力边界。Code Llama 在通用语言模型基础上针对代码数据进行训练，覆盖代码补全、代码生成和自然语言到代码等任务\cite{Roziere2023CodeLlama}。DeepSeek-Coder 则强调高质量代码语料和代码智能能力的系统构建，在多种代码任务上展现出较强表现\cite{Guo2024DeepSeekCoder}。这类模型说明，代码大模型已经具备较强的通用编程能力，为下游任务微调和强化学习优化提供了良好基础。

与此同时，研究者也逐渐意识到，单纯生成代码并不意味着代码真正正确。Liu 等人提出 EvalPlus，指出已有 HumanEval 等基准中的测试用例数量有限，可能高估模型生成代码的真实正确率，并通过扩展测试用例更严格地评估模型功能正确性\cite{Liu2023EvalPlus}。Coder-Reviewer reranking 等工作则尝试通过额外的审查模型或重排序机制提升代码生成结果质量\cite{Zhang2022CoderReviewer}。这些研究共同表明，代码生成模型的评价和优化不能只关注文本相似性或模型似然，而应更多关注执行结果和功能正确性。

\subsection{监督微调在代码生成中的应用}

监督微调是将预训练模型适配到特定任务和数据格式的常用方法。在代码生成任务中，监督微调通常将自然语言问题描述、函数签名或上下文代码作为输入，将参考实现作为输出，使模型学习特定数据集中的解题风格和输出格式。对于小规模实验而言，监督微调具有实现简单、训练稳定、结果可解释等优点，因此常被用作强化学习前的热启动阶段。

然而，监督微调也存在目标错位问题。模型在训练时优化的是参考代码的 token 级似然，而测试时评价的是生成代码是否能通过单元测试。由于同一编程问题通常存在多种等价实现，模型不必生成与参考答案完全相同的代码才能正确；反过来，即使生成代码与参考答案在局部 token 上相似，也可能因为边界条件、返回类型或入口函数错误而无法通过测试。因此，监督微调能够提升模型对任务格式的适应能力，却未必稳定提升功能正确性。

此外，小规模监督微调还可能对强基座模型已有能力产生影响。一方面，领域数据能够强化模型在目标任务上的输出模式；另一方面，数据规模过小或分布较窄时，模型也可能过度拟合特定数据形式，从而削弱原本较强的通用代码能力。近期关于领域监督微调的研究也指出，微调对通用能力的影响取决于数据质量、训练设置和任务分布，并不是简单的单向提升或损伤\cite{Lin2025SFTGeneral}。因此，在本文实验中，SFT 不仅作为独立基线，也作为 PPO 训练的初始策略，用于进一步观察执行反馈能否修复 SFT 后的性能下降。

\subsection{强化学习与执行反馈代码生成研究现状}

强化学习在语言模型后训练中的应用主要源于模型输出质量难以完全通过监督数据刻画的问题。InstructGPT 通过人类反馈训练奖励模型，并利用强化学习优化模型输出，使模型更好地遵循用户指令\cite{Ouyang2022InstructGPT}。PPO 作为一种稳定的策略梯度算法，通过限制新旧策略之间的更新幅度，缓解强化学习训练中的策略崩塌问题\cite{Schulman2017PPO}。在语言模型场景中，PPO 常用于根据奖励信号对生成策略进行进一步优化。

与普通自然语言任务相比，代码任务具有更明确的外部反馈来源。CodeRL 将单元测试结果引入代码生成训练过程，通过深度强化学习提升模型在代码生成任务中的表现\cite{Le2022CodeRL}。PPOCoder 进一步提出将预训练程序语言模型与 PPO 结合，利用代码执行结果和结构对齐信息优化代码生成过程，强调代码的可编译性、语法正确性和功能正确性\cite{Shojaee2023PPOCoder}。COMPCODER 则利用编译器反馈提升神经代码生成结果的可编译性，说明编译错误和执行错误均可以作为有效反馈信号\cite{Wang2022CompCoder}。

除强化学习方法外，也有研究从推理和自我修正角度提升代码生成质量。例如，Chain-of-Thought 提示表明大模型可以通过显式中间推理步骤提升复杂问题求解能力\cite{Wei2022CoT}；Self-Debug 相关工作尝试让模型根据执行反馈自我定位和修复错误\cite{Chen2023SelfDebug}。DPO 等偏好优化方法则提供了不显式训练奖励模型的后训练思路\cite{Rafailov2023DPO}。这些工作说明，后训练阶段的反馈利用方式正在从单一监督学习向多种反馈信号和优化目标发展。

总体来看，国内外相关研究已经证明执行反馈对代码生成任务具有重要价值，但仍存在若干值得进一步研究的问题。首先，如何在小规模数据条件下设计稳定有效的奖励函数仍不明确；其次，强化学习更新可能提升当前验证集表现，却不一定带来跨数据集泛化提升；再次，更复杂的结构奖励是否一定优于简单执行奖励，也需要通过消融实验验证。本文正是在这一背景下，围绕 MBPP 和 HumanEval 数据集构建小规模但完整的 SFT+PPO 实验闭环，重点分析执行反馈奖励在代码生成优化中的实际效果。

\section{本文主要工作}

本文围绕基于执行反馈的强化学习代码生成方法，完成了从数据整理、模型训练到实验分析的完整流程。主要工作如下。

\begin{enumerate}
  \item 构建了统一格式的代码生成实验数据。本文基于 MBPP 和 HumanEval 数据集，将原始数据转换为 master、SFT 和 RL 三类格式，并统一字段为任务编号、问题描述、入口函数、参考答案和测试代码等内容。在此基础上，本文进一步进行了本地 runtime 复核，确认标准答案和测试代码可执行，为后续执行反馈训练提供可靠数据基础。
  \item 实现了基于 Qwen2.5-Coder-1.5B 的监督微调流程。本文使用 MBPP 训练数据构造 prompt-response 样本，采用 LoRA 方式进行参数高效微调，获得 SFT baseline。该阶段为后续 PPO 训练提供初始策略，也用于分析小样本 SFT 对强代码基座模型能力的影响。
  \item 构建了代码生成、执行和打分闭环。本文实现了模型生成代码后的清洗、入口函数对齐、本地执行、单元测试统计和 reward 计算流程，支持在无外部沙箱的服务器环境中完成生成代码的自动评测。该流程将不可微的执行结果转化为强化学习可使用的奖励信号。
  \item 设计并比较了多种执行反馈奖励函数。本文从原始测试通过率奖励出发，进一步设计 reward\_v2，将代码可执行、测试通过比例和全部测试通过奖励组合为更密集的反馈信号。同时，本文参考结构奖励思想设计 reward\_v3 和 reward\_v3b，引入 AST 相似度和函数签名匹配，用于分析复杂结构奖励在小数据条件下的实际效果。
  \item 实现并验证了单步 PPO 优化流程。本文在 SFT 模型基础上收集 rollout，计算 reward 和 advantage，并通过 clipped policy loss、value loss 和 KL penalty 更新 LoRA adapter。实验中进一步比较学习率、KL 系数、rollout 数量和 reward 设计对结果的影响，形成较完整的 PPO 调参消融分析。
  \item 完成了验证集、测试集和案例层面的实验分析。本文比较 Base、SFT 和 PPO 在 MBPP 验证集与 HumanEval 测试集上的表现，并分析 PPO 在部分样例中修复 SFT 错误、在部分样例中产生退化的原因，从而更全面地说明方法的有效性和局限性。
\end{enumerate}

\section{本文组织架构}

全文共分为五章，各章安排如下。

第一章为绪论。该章首先介绍代码生成任务的发展背景和研究意义，然后从代码生成模型、监督微调和执行反馈强化学习三个方面综述国内外研究现状，最后概括本文的主要工作和章节安排。

第二章为相关技术与理论基础。该章将介绍代码生成任务、预训练代码大模型、监督微调、强化学习、PPO 算法、执行反馈优化方法以及常用代码生成数据集和评价指标，为后续方法设计奠定理论基础。

第三章为基于执行反馈的强化学习代码生成方法。该章将详细说明本文的数据处理流程、SFT 训练数据构造、代码生成与执行反馈流程、奖励函数设计、PPO 训练流程和系统实现模块。

第四章为实验设计与结果分析。该章将介绍实验环境、数据集统计、SFT 实验、PPO 主实验、超参数消融实验、奖励函数消融实验、HumanEval 泛化分析和典型案例分析，并总结实验结论。

第五章为全文总结与展望。该章将总结本文研究所得主要结论，分析当前工作在数据规模、训练框架、奖励设计和泛化能力方面的局限，并提出后续研究方向。

\section{本章小结}

本章围绕基于执行反馈的强化学习代码生成研究，首先介绍了代码大模型和代码生成任务的发展背景，指出代码生成任务相比普通文本生成更强调可执行性和功能正确性。随后，本章从代码生成模型、监督微调和执行反馈强化学习三个方面梳理了相关研究现状，说明现有方法虽然已经取得显著进展，但仍存在训练目标与功能正确性不一致、执行反馈利用不足以及小规模数据下强化学习稳定性不确定等问题。在此基础上，本章概括了本文的主要研究工作和全文组织结构。后续章节将进一步介绍相关理论基础、方法实现和实验结果分析。
